\documentclass[10pt,a4paper]{amsart}
\usepackage[margin=19mm]{geometry}
\usepackage{amsmath,amssymb,mathtools}
\usepackage[scr=rsfs,cal=euler]{mathalfa}
\usepackage{xfrac}
\usepackage[unicode,hidelinks,bookmarks=false]{hyperref}
\newcommand{\ie}{i.e.\ }
\newcommand{\cf}{cf.\ }
\newcommand{\resp}{respectively\ }
\newcommand{\Subsection}{\subsection}
\providecommand{\nobreakdash}{\nobreak\hbox{-}}
\setlength{\emergencystretch}{2em}
\multlinegap0pt
\usepackage{amssymb}
\newtheorem{thm}{Theorem}
\title{Alexander polynomials of ribbon knots and virtual knots}
\author{Sheng Bai}
\date{Source: arXiv:2103.07128v3 (CC BY 4.0)}
\begin{document}
\maketitle

\noindent\emph{Source and licence.} Alexander polynomials of ribbon knots and virtual knots. Version arXiv:2103.07128v3 (CC BY 4.0), distributed by arXiv under the Creative Commons Attribution 4.0 International licence, http://creativecommons.org/licenses/by/4.0/, as recorded in the arXiv licence metadata for this exact version and shown on the abstract page https://arxiv.org/abs/2103.07128v3. Full licence notice: \texttt{licenses/arxiv-2103.07128-CC-BY-4.0.txt}.\par\smallskip

\noindent\textit{Editorial scope.} Counted unit: the article's thm thm:algo2 (source0.tex lines 673--675). The article's complete original proof of it is delivered with it (source0.tex lines 677--749, one \texttt{proof} environment of 73 lines). No mathematics is written, repaired or re-derived: the statement and the proof are the article's own lines, in the article's own notation, and no retained line is substituted or re-flowed.  Retained uncounted as the source-local context the proof needs: lines 163--163; lines 185--201; lines 210--226; lines 648--653.  One declared adaptation: exactly 1 ancillary strings inside the retained spans are omitted (image-inclusion commands and, where the source repeats a label, the repeated label definitions) and no other line differs from the source; the figure environments, with their placement, captions and labels, are kept, and the package records the omitted strings in fidelity receipts bound to the affected bodies.  No retained line contains a citation, so the wrapper carries no bibliography and no bibliography entry.  The retained lines contain the article's own diagram code, which the wrapper typesets with the standard tikz packages; no image file is delivered and the package declares no asset dependency.  Editorial formatting only, in addition: an AMS \texttt{amsart} preamble that restates the article's own declarations and macros used by the retained lines, namely the environments \texttt{thm} on separate counters, together with the standard packages those lines need.  The retained environments print in this document as Theorem 1 in order of appearance, whereas the article prints the same items with the numbering of its own full document.  The complete native source of the article as distributed in the arXiv e-print for this exact version is bundled unchanged as \texttt{sources/109076/source0.tex}.
\subsection{Ribbon graph, half Alexander polynomial and our basic theorem}\label{subsect:halfalexpoly}

Now we define a \textit{ribbon matrix} $\rho_{g \times g}$ in terms of the ribbon graph. For each $i = 1, \dots, g$, subdivide $E_i$. Denote the new vertex by $v_{E_i}$ and the obtained directed tree by $T_{E_i}$. Let $P_i$ be the unique path in $T_{E_i}$ from $v_{E_i}$ to $S(E_i)$. Let
\begin{align}
    \rho_{ii} = 
\begin{cases}
\frac{1}{2}, & \text{if the first edge of } P_i \text{ is a forward arc;} \\
-\frac{1}{2}, & \text{if the first edge of } P_i \text{ is a reverse arc.}
\end{cases}
\end{align}
\begin{align}
\rho_{ji} = 
\begin{cases}
1, & \text{if } E_j \text{ is a forward arc of } P_i; \\
-1, & \text{if } E_j \text{ is a reverse arc of } P_i; \\
0, & \text{if } E_j \notin P_i.
\end{cases}
\quad \forall j \neq i. \nonumber
\end{align}\label{for:ribbon matrix rho}

Set matrix $R(t)_{g \times g} = (t - 1)\rho - \frac{1}{2}(t + 1)I$. That is,
\begin{align}
    R_{ii}(t) = 
\begin{cases}
-1, & \text{if the first edge of } P_i \text{ is a forward arc;} \\
-t, & \text{if the first edge of } P_i \text{ is a reverse arc.}
\end{cases}
\end{align}
\begin{align}
    R_{ji}(t) = 
\begin{cases}
t - 1, & \text{if } E_j \text{ is a forward arc of } P_i; \\
1 - t, & \text{if } E_j \text{ is a reverse arc of } P_i; \\
0, & \text{if } E_j \notin P_i.
\end{cases}
\quad \forall j \neq i. \nonumber
\end{align}\label{for:matrix R(t)}

\begin{figure}
    \centering
    
    \caption{Contracted ribbon graph.}
    \label{fig:contracted ribbon graph}
\end{figure}

\begin{thm}\label{thm:algo2}
The half Alexander polynomial $A_{R}(t)=\pm |( {\textbf{R}}_{ji} (t) )| $.
\end{thm}

\begin{proof}
We begin by considering ribbon graph $(T,S)$. After relabelling, for simplification of notations, we may assume $\hat{\gamma}_{1},\hat{\gamma}_{2},...,\hat{\gamma}_{\omega}$ are successive arcs in the ribbon diagram as in Fig. \ref{fig:contracted ribbon graph}(1) so that the path $E_{1}E_{2}...E_{\omega}$ in $T$ corresponds to an edge $E$ in $T^{\omega}$ with $\omega=\omega(E)$. We may assume $\hat{\gamma}_{2},...,\hat{\gamma}_{\omega}$ are in the path $P_{1}$ as defined before (\ref{for:ribbon matrix rho}) in Subsection \ref{subsect:halfalexpoly}.

Suppose $\hat{\gamma}_{1}$ is a forward arc of $P_{1}$. By (\ref{for:matrix R(t)}), the minor

\begin{equation}
(R_{ji}(t))_{1,...,\omega;1,...,\omega}=\left(\begin{array}[]{ccccc}
-1 & & & & \\
t-1 & -1 & & & \\
t-1 & t-1 & -1 & & \\
\vdots & \vdots & \vdots & \ddots & \\
t-1 & t-1 & t-1 & \cdots & -1
\end{array}\right).
\end{equation}

It is easy to see from ribbon graph that

\begin{align}
R_{ji}(t) &=R_{j,\omega}(t), &\quad \forall j\neq 1,2,...,\omega, \quad \forall i=1,2,...,\omega; \nonumber \\
R_{ji}(t) &=R_{j,\omega+1}(t), &\quad \forall j=1,2,...,\omega, \quad \forall i\neq 1,2,...,\omega. \label{for:Rijt 1tow rows}
\end{align}

For $i=1,2,...,\omega-1$, subtract the $(i+1)$th column from the $i$th column, which means choosing $[f_{1}]-[f_{2}],...,[f_{\omega-1}]-[f_{\omega}]$ instead of $[f_{1}],...[f_{\omega-1}]$ in the basis. Then the minor becomes

\begin{equation}
(R_{ji}(t))_{1,...,\omega;1,...,\omega}=\left(\begin{array}[]{cccccc}
-1 & & & & \\
t & -1 & & & \\
0 & t & -1 & & \\
\vdots & \vdots & \vdots & \ddots & \\
0 & 0 & 0 & & -1 \\
0 & 0 & 0 & \cdots & t & -1
\end{array}\right)
\end{equation}
and now

\begin{equation}
R_{ji}(t)=0, \quad \forall j\neq 1,2,...,\omega, \quad \forall i=1,2,...,\omega-1. \nonumber
\end{equation}

Multiply the $1$st, $2$nd ,..., $(\omega-1)$th rows by $t^{\omega-1},t^{\omega-2},...,t$ respectively and add together to the $\omega$th row. Then the new $R_{\omega,\omega-1}(t)$ is $0$. By (\ref{for:Rijt 1tow rows}), for $i\neq 1,2,...,\omega$, the new $R_{\omega,i}(t)$ is $(1+t+\cdots+t^{\omega-1})$ times the old $R_{\omega,i}(t)$. Now we see

\begin{equation}
|\left(R_{ji}(t)\right)|=(1+t+\cdots+t^{\omega-1})\left|\begin{array}[]{cccccc}
-1 & & & & 0 & \\
t & -1 & & & \vdots & *\\
\vdots & \vdots & \ddots & & \vdots & \\
0 & 0 & \cdots & -1 & 0 &  \\
0 & 0 & \cdots & 0  & \frac{-1}{1+t+\cdots+t^{\omega-1}}  & * \\
O &   &        & *  & * & *
\end{array}\right|
\end{equation}
where $\frac{-1}{1+t+\cdots+t^{\omega-1}}$ is the $(\omega,\omega)-$entry. Thus

\begin{equation}
|\left(R_{ji}(t)\right)|=(-1)^{\omega-1}(1+t+\cdots+t^{\omega-1})|\left(R^{*}_{ji}(t)\right)|,
\end{equation}
where $(R^{*}_{ji}(t))$ is the minor of the original $(R_{ji}(t))$ deleting the $1$st, $2$nd ,..., $(\omega-1)$th rows and the $1$st, $2$nd ,..., $(\omega-1)$th columns and replacing $-1$ by $\frac{-1}{1+t+\cdots+t^{\omega-1}}$ at the original $(\omega,\omega)-$entry. Notice that $(1-t)(1+t+\cdots+t^{\omega-1})=1-t^{\omega}$. As a result,

\begin{equation}
|\left(R_{ji}(t)\right)|=(-1)^{\omega(e)-1}|\left(R_{ji}^{**}(t)\right)|, \label{for:Rijt replace t to tw}
\end{equation}
where $\left(R_{ji}^{**}(t)\right)$ is the minor of the original $\left(R_{ji}(t)\right)$ deleting the $1$st, $2$nd ,..., $(\omega-1)$th rows and the $1$st, $2$nd ,..., $(\omega-1)$th columns and replacing $t$ by $t^{\omega}$ in the original $\omega$th column.

If $\hat{\gamma}_{1}$ is a reverse arc of $P_{1}$, A similar argument shows that (\ref{for:Rijt replace t to tw}) still holds.

Do it inductively for path on $T$ corresponding to each weighted edge in $T^{\omega}$, we finally have

\begin{equation}
|\left(R_{ji}(t)\right)|=(-1)^{g-|T^{\omega}|}|\left(R_{ji}(t)\right)|, \nonumber
\end{equation}
where $g$ is the number of edges in $T$ and $|T^{\omega}|$ denotes the number of edges in $T^{\omega}$.
\end{proof}

\end{document}
